% !TeX root = ../main.tex
% Add the above to each chapter to make compiling the PDF easier in some editors.

\chapter{Conclusion}\label{chapter:conclusion}

This thesis investigated how accurately a decision tree model can predict query runtimes on an iterator-based database system such as PostgreSQL.
Furthermore, it analyzed which of T3's key concepts -- pipeline-based representation and tuple-centric prediction targets -- transfer effectively, and how the result compares to existing learned cost models.

We adapted T3 for PostgreSQL by converting parsed ZeroShot execution plans into T3's pipeline representation and constructing a PostgreSQL-native feature vector.
The feature vector includes PostgreSQL-specific signals such as planner cost estimates and planned parallelism alongside operator-level features derived from T3's Umbra design.
We trained a LightGBM model with a leave-one-database-out split, mirroring both T3's and ZeroShot's generalization evaluation.

The model generalizes reliably across unseen PostgreSQL database instances within the ZeroShot workload, with a median Q-Error below 1.33 across all 20 ZeroShot databases.
On the Join Order Benchmark, it clearly outperforms the T3-based PostgreSQL implementation of Wehrstein et al.~\parencite{jobcomplex}.
Hence, this result suggests that a PostgreSQL-native feature vector is beneficial.
Among all learned cost models evaluated on JOB, our model ranks 4th with a median Q-Error of 2.42.
This is well ahead of most learned approaches and far better than PostgreSQL's built-in estimator at over 1000.
However, it falls behind SOTA neural network models such as ZeroShot (1.58) and DACE (1.91), and behind T3 (Umbra).

Regarding T3's key concepts, the pipeline-based representation transfers well.
It provides the structural basis for feature extraction and produces a compact, informative feature vector.

However, tuple-centric prediction targets do not transfer.
PostgreSQL's iterator model causes per-node timing intervals to overlap across pipeline boundaries, so true exclusive execution times per pipeline are unrecoverable.
Per-query prediction therefore works best -- the opposite of T3's finding for Umbra.

Finally, a promising path to closing the gap to T3 (Umbra) and other SOTA models is a dedicated PostgreSQL benchmark covering a wider range of query types and runtimes.
With such a benchmark, the per-tuple approach could be re-evaluated, especially after addressing the pipeline timing issue.
A purely query-level feature vector -- bypassing pipeline decomposition -- could also be investigated.
Moreover, prediction latency should be investigated using native code compilation with the \texttt{lleaves} framework.

Overall, this work shows that T3's approach is viable beyond Umbra.
Our results establish a competitive PostgreSQL baseline that future work can build on.
